\begin{papercard}{page2021}
\cf{Stages} \stagetag{R} (reporting items and flow diagram), \stagetag{A} (reporting how apps or sources were identified, screened and excluded), \stagetag{D} (data items and availability of materials). PRISMA 2020 statement (J.\ Clin.\ Epidemiol.\ 134:178--189; here 13 PDF pages including references): the 27-item reporting checklist for systematic reviews, a separate 12-item abstract checklist and a flow-diagram template. It is a \emph{reporting} guideline for systematic reviews, not a method; the thesis is not a systematic review, so the card lists items that can be adapted to describing app and literature selection.

\cf{Purpose} PRISMA 2009 was designed to make authors report ``why the review was done, what the authors did, and what they found''; the 2020 version replaces it to reflect new methods for identifying, selecting, appraising and synthesising studies \pc{page2021}{3}{was designed to help systematic reviewers transparently report why the review was done, what the authors did, and what they found}. Advances prompting the update: text mining and machine learning for study identification, synthesis without meta-analysis, new risk-of-bias tools, certainty-of-evidence language, registration of protocols, preprints and repositories.

\cf{Development of PRISMA 2020} Items often reported incompletely in earlier reviews were identified from meta-research; 60 documents with reporting guidance were reviewed; an online survey of systematic-review methodologists and editors (110 of 220 invited responded); a 21-member, 2-day in-person meeting in Edinburgh in September 2018; drafts in 2019--2020 and a layout/terminology survey (15 of 22 responded) \pc{page2021}{4}{Systematic review methodologists and journal editors were invited to complete the online survey (110 of 220 invited responded)}.

\cf{Scope and limits of the guideline} Designed primarily for reviews of health interventions, but applicable to reviews of other interventions and, in many items, to reviews with other objectives (etiology, prevalence, prognosis); for qualitative synthesis consult other guidelines \pc{page2021}{4}{The PRISMA 2020 statement has been designed primarily for systematic reviews of studies that evaluate the effects of health interventions}. ``PRISMA 2020 is not intended to guide systematic review conduct'' and ``should not be used to assess the conduct or methodological quality of systematic reviews''; protocols have PRISMA-P; extensions exist for network meta-analysis, IPD, harms, diagnostic accuracy and scoping reviews (PRISMA-ScR) \pc{page2021}{5}{PRISMA 2020 should not be used to assess the conduct or methodological quality of systematic reviews}.

\cf{Structure} Seven sections, 27 items with sub-items \pc{page2021}{5}{The PRISMA 2020 checklist includes seven sections with 27 items, some of which include sub-items}; an explanation-and-elaboration paper gives bullet-point ``elements''; the location of items in the report is not prescriptive, ``the guiding principle is to ensure the information is reported'' \pc{page2021}{6}{the suggested location should not be seen as prescriptive; the guiding principle is to ensure the information is reported}; supplementary files may be placed in an open repository (OSF, Dryad, figshare).

\cf{Main changes from 2009 (Box 2)} Abstract checklist included; protocol and registration moved to ``Other information'' with amendments (24c); full search strategies for \emph{all} databases and sources \pc{page2021}{5}{Modification of the ‘Search’ item to recommend authors present full search strategies for all databases, registers and websites searched, not just at least one database}; selection item stresses how many reviewers screened each record and report, independence, and automation tools; data items define outcomes (10a); synthesis item split into six sub-items (13a--f); list of excluded studies with reasons (16b); certainty-of-evidence items (15, 22); competing interests (26); availability of data, code and materials (27) \pc{page2021}{6}{Addition of a new item recommending authors indi- cate whether data, analytic code and other materials used in the review are publicly available and if so, where they can be found (see item \#27)}.

\cf{The checklist (Table 1)} \emph{Title} 1 identify as a systematic review. \emph{Abstract} 2. \emph{Introduction}: 3 rationale; 4 objectives. \emph{Methods}: 5 eligibility criteria and grouping for synthesis; 6 information sources with last-search dates \pc{page2021}{7}{Specify all databases, registers, websites, organisations, reference lists and other sources searched or consulted to identify studies. Specify the date when each source was last searched or consulted}; 7 full search strategies including filters and limits; 8 selection process (number of reviewers, independence, automation) \pc{page2021}{7}{including how many reviewers screened each record and each report retrieved, whether they worked independently}; 9 data collection process; 10a--b data items (outcomes and other variables); 11 risk-of-bias assessment; 12 effect measures; 13a--f synthesis methods (eligibility for each synthesis, data preparation, tabulation/visual display, synthesis method with rationale, heterogeneity exploration, sensitivity analyses); 14 reporting-bias assessment; 15 certainty assessment. \emph{Results}: 16a search and selection results with flow diagram; 16b excluded studies with reasons \pc{page2021}{8}{Cite studies that might appear to meet the inclusion criteria, but which were excluded, and explain why they were excluded}; 17 study characteristics; 18 risk of bias per study; 19 results of individual studies; 20a--d synthesis results; 21 reporting biases; 22 certainty. \emph{Discussion}: 23a--d interpretation, limitations of evidence, limitations of the review process, implications. \emph{Other}: 24a--c registration, protocol access, amendments \pc{page2021}{8}{Provide registration information for the review, including register name and registration number, or state that the review was not registered}; 25 support; 26 competing interests; 27 availability of data, code and materials \pc{page2021}{8}{Report which of the following are publicly available and where they can be found}.

\cf{Abstract checklist and flow diagram} Twelve items: title, objectives, eligibility, information sources and dates, risk of bias, synthesis methods, included studies, synthesis results, limitations of evidence, interpretation, funding, registration; retains the 2013 PRISMA-for-Abstracts items plus a synthesis-method item \pc{page2021}{9}{This abstract checklist retains the same items as those included in the PRISMA for Abstracts statement published in 2013}. Figure 1 is a flow template (identification from databases and registers/other methods, duplicates and automation removals, screening, reports sought and assessed, exclusions with reasons, studies included); boxes in grey are optional \pc{page2021}{10}{PRISMA 2020 flow diagram template for systematic reviews}.

\cf{Discussion by the authors} Complete reporting lets readers judge the methods and trustworthiness and supports replication and updating; strategies to raise adherence (education, endorsement by journals, reviewers checking, authors indicating where items are reported, writing tools) have rarely been evaluated: of 31 proposed interventions, effects of only 11 were evaluated, mostly observational and at high risk of bias; they suggest think-aloud and reliability studies to find items interpreted inconsistently.

\cf{Limits (card)} Health-intervention orientation; reporting only (a well-reported review can be poorly conducted); no guidance on qualitative or app-store searches; checklist adherence is not a quality score; the paper itself notes items might be interpreted inconsistently and that reliability of item interpretation has not been tested.

\cf{Use for this thesis} (1) The thesis' app-selection and literature-search sections can adopt the reporting logic without claiming to be a systematic review: eligibility criteria (5), sources and dates (6), the exact search/filter settings (7), who screened and how independently (8), a flow of numbers from identified to included with reasons for exclusion (16a--b), data items and their definitions (10), and an availability statement for codebooks, extracted data and code (27). (2) A statement of any deviation from a pre-specified plan (24c) fits the Decision Log. (3) For the scoping-review character of parts of the work cite PRISMA-ScR (Tricco et al.) and PRISMA-S for searches (Rethlefsen et al.) alongside this paper. (4) Do not claim that the thesis ``follows PRISMA'' unless the checklist has been completed and attached; say instead which items were adapted.
\end{papercard}
